\chapter*{Tóm tắt đồ án}
\addcontentsline{toc}{chapter}{Tóm tắt đồ án}

Việt Nam là quốc gia có địa hình hẹp dài, khí hậu nhiệt đới gió mùa phân hóa phức tạp và thường xuyên chịu ảnh hưởng bởi các hình thái thời tiết cực đoan như bão lũ, mưa lớn, sương mù dày đặc và dông sét. Các yếu tố này gây ra nguy cơ mất an toàn giao thông nghiêm trọng cho người điều khiển phương tiện, đặc biệt là người đi xe máy -- phương tiện di chuyển phổ biến nhất tại Việt Nam.

Đồ án này nghiên cứu, thiết kế và hiện thực hóa \textbf{“Hệ thống thông tin thời tiết thông minh và tư vấn lộ trình di chuyển tại Việt Nam tích hợp AI”}. Hệ thống giải quyết trọn vẹn các thách thức cốt lõi:
\begin{enumerate}
    \item \textit{Thu thập và chuẩn hóa dữ liệu khí tượng}: Tự động hóa tiến trình thu thập dữ liệu thời tiết 63 tỉnh thành thời gian thực, lưu trữ trên cơ sở dữ liệu không gian PostGIS cùng cơ chế kiểm soát nghiêm ngặt tính xác thực và độ tươi mới của dữ liệu.
    \item \textit{Động cơ đánh giá rủi ro thời tiết đa chiều}: Xây dựng ma trận phân tích rủi ro dựa trên 5 khía cạnh: cường độ mưa, gió giật, tầm nhìn, nhiệt độ cực đoan và độ trơn trượt đường sá, cá nhân hóa theo từng loại phương tiện di chuyển (ô tô và xe máy).
    \item \textit{Thuật toán lập lộ trình cân bằng đa mục tiêu}: Tích hợp dữ liệu thời tiết vào dịch vụ định tuyến địa lý để đề xuất đồng thời 3 hành lang di chuyển: Nhanh nhất, Ngắn nhất và An toàn thời tiết nhất, hỗ trợ khoảng cách liên tỉnh lên tới 2.000~km.
    \item \textit{Trợ lý ảo AI có kiểm chứng dữ liệu (AI Grounding)}: Ứng dụng mô hình ngôn ngữ lớn (Google Gemini) kết hợp cơ chế Function Calling nội bộ nhằm đảm bảo câu trả lời luôn được trích dẫn dữ liệu quan trắc thật, loại bỏ hoàn toàn hiện tượng ảo tưởng thông tin.
    \item \textit{Kiến trúc phân tán an toàn và kiểm soát hạn ngạch}: Hệ thống được triển khai trên nền tảng Docker với mô hình phân vùng mạng an toàn, kiểm soát hạn ngạch nguyên tử trên Redis theo múi giờ Việt Nam và tích hợp cổng thanh toán trực tuyến.
\end{enumerate}

\textbf{Từ khóa:} Khí tượng giao thông, Định tuyến thời tiết, PostGIS, Đánh giá rủi ro, Mô hình ngôn ngữ lớn, AI Grounding, Fastify, MapLibre GL.

\vspace{10mm}
\begin{center}
\textbf{\large ABSTRACT}
\end{center}

Vietnam possesses a long and narrow geography characterized by complex tropical monsoon climates and frequent severe weather phenomena such as typhoons, torrential rains, dense fog, and thunderstorms. These meteorological factors pose critical travel safety risks, particularly for motorcyclists, who constitute the predominant proportion of commuters in the country.

This graduation thesis presents the research, design, and implementation of the \textbf{“Smart Vietnam Weather Information and Safe Travel Route Planning System with AI Grounding”}. The system addresses these key engineering challenges:
\begin{enumerate}
    \item \textit{Meteorological Ingestion and Spatial Normalization}: Automated real-time harvesting across 63 provinces, persisted within a PostGIS database with continuous validation of data authenticity, integrity, and temporal freshness.
    \item \textit{Config-Driven Multi-Dimensional Risk Engine}: Evaluation across precipitation intensity, wind gusts, surface visibility, extreme temperatures, and road surface conditions, customized by vehicle profile.
    \item \textit{Multi-Alternative Weather-Aware Route Optimization}: Concurrent calculation of Fastest, Shortest, and Weather-Safer routes for inter-provincial corridors spanning up to 2,000~km.
    \item \textit{Grounded AI Travel Assistant}: Integration with Google Gemini LLM via function calling over internal tool declarations, ensuring responses are grounded with verifiable citations without hallucinations.
    \item \textit{Containerized Architecture and Quota Governance}: Segmented Docker networking, atomic Redis quota counters aligned with Asia/Ho Chi Minh midnight resets, and integrated digital payments.
\end{enumerate}

\textbf{Keywords:} Weather-Aware Routing, Travel Risk Assessment, PostGIS, LLM Grounding, Fastify, MapLibre GL, Vietnam Climate.
